% Separate post-CP29 candidate; not part of the frozen paper or current gate.
\subsection{Two opposite actual quotients and the first wild matrix layer}

Let \(K\) be complete discretely valued of mixed characteristic
\((0,p)\), with algebraically closed residue field,
\(p\ge5\), \(e_K=e=p(p-1)\), and \(\mu_p\subset K\).
Put \(q=p^2\), \(N=q-1\), \(s=(p+1)e=pN\), and use
the common tame good field \(L/K\) of degree \(N\).
All groups considered are height-two Barsotti--Tate groups
acquiring one-dimensional supersingular formal good reduction.
All torsion comparisons and quotients below are actual ones.

\begin{lemma}[The tame second-level point-field test at this index]
\label{lem:B-RD83-tame}
\status{Candidate; full proof supplied, independent audits pending.}
Suppose two good formal height-two groups over \(L\) have
two-slope first-level polygons with parameters \(B_i=pb_i\),
where \(p\nmid b_i\), and have isomorphic actual
\(p^2\)-torsion. If both first-level modules are
\(L\)-rational, then \(b_1=b_2\).
\end{lemma}
\begin{proof}
This is the common-point-field proof underlying B69, written
again without its prime-to-\(p\) absolute-index assumption.
Choose corresponding residual vectors outside the two
canonical high lines, and corresponding primitive
second-level lifts \(P_i\). The actual inverse Newton
polygon gives \(v_L(T_i(P_i))=b_i/q\). The torsion
isomorphism identifies their stabilizers, hence their
point fields \(E=L(P_1)=L(P_2)\).

Because the first-level modules are \(L\)-rational and
\(pP_i\) is rational, each point field is Galois, with
its group a subgroup of the first-level translation
group. Indeed every conjugate is \(P_i+Q\) for an
\(L\)-rational \(Q\), so it is defined over the same
point field. Thus \([E:L]\le q\). The reduced denominator
of \(b_i/q\) is \(q\), so \([E:L]=q\), and every
first-level translation occurs.

Choose a translation outside both high lines. With the
integral valuation on \(E\), the two primitive coordinates
have valuations \(b_i\), and their low displacements have
valuations \(qb_i\). The exact prime-valuation equality
in Lemma~\ref{lem:B-RD60-break} gives its common lower
break as \((q-1)b_1=(q-1)b_2\). This proves the assertion.
\end{proof}

\begin{lemma}[Both exchanged-line quotients have residual wild action]
\label{lem:B-RD83-quotients}
\status{Candidate; full proof supplied, independent audits pending.}
Suppose two groups with isomorphic actual \(p^3\)-torsion
have tame reducible residual inertia and two-slope parameters
\[
                 B_i=p^2k_i,\qquad1\le k_i\le p-2.
                                                        \tag{B-RD83.1}
\]
Their canonical high lines are assumed to be distinct under
the comparison. Quotienting both groups by either of the two
common stable lines gives an actual \(p^2\)-torsion comparison
with nontrivial residual wild inertia over \(L\).
\end{lemma}
\begin{proof}
The original first-level modules are \(L\)-rational. Each
of the two lines is stable under the original inertia action:
one is the high line of the first model and the other is
the high line of the second. Form their actual schematic
closures in the good models and the corresponding finite-flat
quotients. The generic quotients descend to \(K\) by the
actual semilinear descent argument in B69; no good integral
model over \(K\) is asserted.

Every parameter (B-RD83.1) is greater than \(e\).
The exact high and low quotient parameters are therefore
\[
            B_{i,H}=s-B_i=p(N-pk_i),
       \qquad B_{i,J}=B_i/p=pk_i\quad(J\ne H_i).
                                                        \tag{B-RD83.2}
\]
For clarity, these formulas are the actual Weierstrass
product calculation of Lemma~\ref{lem:B-RD69-quotients},
which does not use \(p\nmid e\). For the high quotient,
each inverse of a high target has valuation \(\alpha/q\)
because \(B>e\); quotienting multiplies this by \(p\),
while the old first-level image has valuation \(B\).
The resulting parameter is \(\alpha=s-B\). For a low
quotient the old high representative has quotient valuation
\(\alpha+(p-1)\beta=s-B/p\); the inverse low points
give the other valuation \(B/p^2\), hence parameter \(B/p\).

The actual identity
\[
 (M/J)[p^2]=\{P\in M[p^3]:p^2P\in J\}/J              \tag{B-RD83.3}
\]
gives the asserted comparisons. For either chosen line,
one model uses a high quotient and the other a low quotient.
Their parameters divided by \(p\) are \(N-pk_i\) and
\(k_j\), both prime to \(p\), and
\[
                  N-pk_i\ge2p-1>p-2\ge k_j.          \tag{B-RD83.4}
\]
The original third-level image over \(L\) is a
\(p\)-group, since its first-level module is rational.
The actual quotient comparison (B-RD83.3) is a subquotient
of that module, so its image over \(L\) is also a
\(p\)-group. If its residual action were tame, it would
therefore be trivial. The preceding lemma would then
equate these unequal integers. Thus its residual wild
action is nontrivial. The actual old first-level image
\(M[p]/J\subset(M/J)[p]\) also supplies an invariant
line, so this is a nontrivial unipotent residual action.
\end{proof}

\begin{proposition}[Opposite quotient wild actions force five matrix dimensions]
\label{prop:B-RD83-Lie}
\status{Candidate; full proof supplied, independent audits pending.}
Assume in addition \(p\equiv1\pmod4\) and that the common
finite determinant is nonsquare. Under the hypotheses of
Lemma~\ref{lem:B-RD83-quotients}, the original actual
third-level wild image over \(L\) has order at least \(p^5\).
\end{proposition}
\begin{proof}
Identify the actual third-level torsion modules by the given
isomorphism, and put the two stable lines at
\(\mathbf F_pe_1\) and \(\mathbf F_pe_2\). The common
residual representation is tame and diagonal over
\(\mathbf F_p\). All its action is killed by \(L\), so
\[
 W=\operatorname{Gal}(L(M[p^3])/L)
        \subset1+pM_2(\mathbf Z/p^3).
\]
Let \(W_1\) and \(W_2\) be its actual first and terminal
matrix spaces, as in (B-RD79.2). The actual matrix identities
give
\[
 W_1\subset W_2,\qquad [W_1,W_1]\subset W_2,
       \qquad |W|=p^{\dim W_1+\dim W_2}.              \tag{B-RD83.5}
\]

We identify the residual actions of the two quotients at
the integral Tate-lattice level. If \(T=T_pM\), then under
the rational identification supplied by the isogeny the
quotient by the first line has lattice
\(T'=T+\mathbf Z_p(e_1/p)\). In the basis
\((e_1/p,e_2)\), a matrix \(\rho=(\rho_{ij})\) becomes
\[
                \begin{pmatrix}
                 \rho_{11}&p\rho_{12}\\
                 \rho_{21}/p&\rho_{22}
                \end{pmatrix}.                       \tag{B-RD83.6}
\]
For \(\rho(g)=1+pA+p^2B\), its residual lower entry is
\(A_{21}\); the other off-diagonal residual entry is zero.
Thus the nontrivial residual wild action of this quotient
forces a nonzero \(21\)-projection of \(W_1\). Quotienting
by the second line gives, in the same way, a nonzero
\(12\)-projection. These are coefficients of the original
actual representation, not independently chosen matrices.

Let a residual tame generator have eigenvalues
\(\alpha,\beta\in\mathbf F_p^\times\). Since
\(\chi_p\equiv1\pmod p\), their product is the residual
finite determinant and is nonsquare. Therefore
\(\rho_0=\alpha/\beta=\alpha^2/(\alpha\beta)\) is
nonsquare as well. When \(p\equiv1\pmod4\), both
\(1\) and \(-1\) are squares. Hence
\(1,\rho_0,\rho_0^{-1}\) are distinct.

Tame conjugation preserves \(W_1\). Its three eigenspaces
in \(M_2(\mathbf F_p)\) are the diagonal plane,
\(\mathbf F_pE_{12}\), and \(\mathbf F_pE_{21}\), with
those distinct weights. The polynomial spectral projectors,
whose denominators are nonzero elements of \(\mathbf F_p\),
preserve \(W_1\). The two nonzero off-diagonal projections
therefore imply \(E_{12},E_{21}\in W_1\). Their commutator
is \(\operatorname{diag}(1,-1)\), independent from both.
Equation (B-RD83.5) now gives
\(\dim W_1\ge2\) and \(\dim W_2\ge3\), proving the bound.
\end{proof}

\begin{corollary}[One lower bound across all residual types at \(p=5\)]
\label{cor:B-RD83}
\status{Candidate; full proof supplied, independent audits pending.}
For \(p\equiv1\pmod4\), a different-special-character
actual \(p^3\)-congruent pair at \(e=p(p-1)\) has
third-level wild image of order at least \(p^5\), regardless
of its residual type. In particular any remaining
\(125\)-torsion ambiguity at \(p=5,e=20\) has wild image
of order at least \(3125\).
\end{corollary}
\begin{proof}
The square finite-determinant case is excluded by B77, so
assume the common finite determinant is nonsquare.
Irreducible residual inertia is covered by B81, and
nontrivial residual wild inertia by B82.

For tame reducible residual inertia a one-slope polygon
would have a leading-coordinate plane with scalar action
over \(\mathbf F_p\), and hence square determinant.
Thus both polygons have two slopes. Their first-level
modules are rational over \(L\), so \(B_i=pb_i\).
If one \(b_i\) is prime to \(p\), its primitive third-level
point has valuation \(b_i/p^4\), forcing a wild field
degree divisible by \(p^4\). Equality of that degree
with \(p^4\) recovers the character by B82. Hence an
unequal-character pair in this case has order at least
\(p^5\).

If both \(b_i\) are divisible by \(p\), write
\(B_i=p^2k_i\) with \(1\le k_i\le p-2\). If the high
lines agree, equality of finite determinants gives
\(c_1\equiv c_2\pmod{p-1}\); equality of their high-line
characters gives \(B_1\equiv B_2\pmod{p-1}\). The stated
range then gives \(k_1=k_2\), and the line character gives
\(c_1=c_2\pmod N\). Otherwise the lines are distinct,
and the proposition applies. These alternatives exhaust
the residual types.
\end{proof}

\begin{remark}[The still-open larger-image comparisons]
\label{rem:B-RD83-scope}
\status{Candidate scope; no exact first-wild index claim.}
No order at least \(p^5\) is excluded or constructed here.
The full first-wild least index remains in \(\{2,3\}\).
The condition \(p\equiv1\pmod4\) is used precisely to
separate the two opposite tame weights. When
\(p\equiv3\pmod4\), a nonsquare ratio can equal \(-1\),
and the two nonzero projections alone do not produce the
two separate matrix units. That case is not included in
the proposition or the final corollary.

\smallskip
\noindent\emph{Exact supplement and provenance.}
The retained standalone program
\path{code/reduction_descent/verify_first_wild_swapped_quotients.py}
checks the two symbolic lattice conjugations and the two
exact residual matrices. For every prime from 5 through 101
it checks all 71,456 opposite-quotient parameter pairs,
36,848 nonsquare residual eigenvalue pairs, 24 literal
commutator and power rows, 24 dimension-bound rows and
74,804 residual-tame parameter branches. The order-two
exceptions at primes congruent to 3 modulo 4 are retained
explicitly, rather than included in the theorem.
The author run passed in \(0.3601258019916713\) seconds,
with no failed execution. The initial and clarified proofs,
every executed program and its output are retained. The
complete internal B69 point-field and actual quotient
proofs were reread and their relevant steps rederived
at the present wild index. No new original-source read
is claimed. These finite checks establish no formal-group
existence and do not replace the uniform proofs.
\end{remark}
